\section{The entire chronological product and its dynamic maximum on open sets}

Throughout this section, the ambient dimension $n$ is positive.

The chronological product is formed from the separately propagated original
members. Its boundary is the full ordered geometric boundary, including every
exceptional member born along the history. The following comparison keeps both
lists in their original roles. It concerns restriction of a given history; it
does not assert decrease or termination of the sharp resolution procedure.

\begin{theorem}[The current product lies in the invariant domain without compactness]
\label{hir:sharp-original-noncompact-current-product-domain}
\uses{hir:sharp-original-chronological-full-product-geometry-order-sum,hir:sharp-original-all-member-actual-finite-chronological-history,hir:geo-transform-wellformed}
\lean{Hironaka.HironakaSharpActualDynamicMaximumOpen.currentProductMarked_isInvDomain_of_pure}
\leanok
Let the original closed smooth carrier have fixed pure dimension
$d+1\leq n$, let the ordered original boundary have simple normal crossings,
and suppose its complement is dense in the carrier. After any finite genuine
geometric test history for the unit template, the actual current marking-one
chronological product belongs to the domain of the canonical invariant in
dimension at most $n+1$. No compactness is required.
\end{theorem}
\begin{proof}
The unit template is well formed, so its transform along the test history
has closed smooth carrier, simple-normal-crossings boundary and the required
transversality. Induct along the history to preserve the original pure
dimension: at the identity it is the hypothesis; at a last blowup the actual
centre lies in the previous carrier, so the strict-transform regularity
theorem preserves dimension $d+1$ at each point of the new carrier. An empty
carrier causes no exception. The true chronological member product has a
coherent ideal and is nonzero along this current carrier by the independently
proved original-member geometry and finite-product theorem. Its marking is
one, and $d+1\leq n+1$. These are exactly the domain requirements. Compactness
does not enter any of these arguments.
\end{proof}

\begin{theorem}[Restriction of the entire current marked product]
\label{hir:sharp-original-current-product-marked-open}
\uses{hir:sharp-actual-relative-finite-product-open,hir:sharp-original-members-whole-chronological-open-naturality,hir:open-restriction-tower}
\lean{Hironaka.HironakaSharpActualDynamicMaximumOpen.currentProductMarked_open}
\leanok
Let $W$ be a closed smooth carrier, $E$ an ordered simple-normal-crossings
boundary, and suppose that $W\setminus |E|$ is dense in $W$. Let $T$ be any
finite geometric test history for the unit marked template of $W$. For an open
set $U$ in the original ambient manifold, pullback to the endpoint over $U$ of
the entire current product marked ideal equals the independently formed
current product for $W\cap U$, the ordered list $E|_U$, and $T|_U$.
\end{theorem}
\begin{proof}
The geometric-transform restriction identity gives the same current carrier
and the same entire ordered boundary on both sides. This includes each born
exceptional member, since the identity concerns the complete marked record of
the unit template. The unit ideal restricts to the unit ideal, so the initial
templates in that identity are exactly those specified in the statement.
For each original position, the chronological restriction theorem identifies
the pullback of its entire current ideal with the independently propagated
ideal on the restricted history. Pullback of the recursively carrier-relative
finite product is the product of these pullbacks. The finite list indexed by
the restricted boundary has the same length as the original list; transport
along this length equality identifies every position without discarding empty
members or identifying coincident members. Thus the entire product ideals
coincide. Both markings equal one. These four equalities prove equality of the
marked records.
\end{proof}

\begin{theorem}[Restriction of every full current invariant value on the actual defect]
\label{hir:sharp-original-current-full-value-open}
\uses{hir:sharp-original-current-product-marked-open,hir:sharp-original-noncompact-current-product-domain,hir:sharp-original-chronological-distinct-defect-open}
\lean{Hironaka.HironakaSharpActualDynamicMaximumOpen.current_full_value_open}
\leanok
In addition, suppose the original carrier is smooth of fixed
positive pure dimension $d+1\leq n$. At every point of the endpoint over $U$ whose image
lies in the actual chronological distinct defect, the full canonical
invariant of the restricted current product equals the full invariant of the
original current product at that image.
\end{theorem}
\begin{proof}
Original geometry and the genuine history place the current product in the
domain of the canonical invariant. Membership in the distinct defect means
that its product order is strictly greater than its number of distinct proper
germ classes, hence at least one. The point therefore lies in the cosupport
of the current marking-one product. Apply the open-restriction law of the full
invariant at the identity history and the endpoint open set. The preceding
whole-record equality identifies its pulled-back marked ideal with the
independently formed restricted current product. This proves equality of the
full values, including all later entries of the invariant.
\end{proof}

\begin{theorem}[A dynamic maximum restricts exactly when its trace is nonempty]
\label{hir:sharp-original-current-dynamic-maximum-open}
\uses{hir:sharp-original-current-full-value-open,hir:sharp-original-chronological-distinct-defect-open}
\lean{Hironaka.HironakaSharpActualDynamicMaximumOpen.current_dynamic_maximum_open_of_meeting}
\leanok
Under the preceding hypotheses, if the endpoint over $U$ meets the maximum
locus of the full current invariant on the actual chronological distinct
defect, then the maximum locus for the independently restricted current data
is exactly its preimage.
\end{theorem}
\begin{proof}
The entire distinct defect restricts to its preimage, retaining the true
quotient of proper germ classes. Choose a point of the original maximum in
the endpoint open. The preceding theorem preserves all full values on that
preimage. Any restricted maximum has value at least that of the chosen point
and at most the original maximum, so it has the original maximum value.
Conversely, any point of the original maximum bounds all restricted values
after the same value identification. These two implications give equality
of the maximum loci.
\end{proof}

\begin{theorem}[Omitting the dynamic maximum strictly lowers every local full value]
\label{hir:sharp-original-current-value-below-omitted-maximum}
\uses{hir:sharp-original-current-full-value-open,hir:sharp-original-chronological-distinct-defect-open}
\lean{Hironaka.HironakaSharpActualDynamicMaximumOpen.current_full_value_lt_of_maximum_trace_empty}
\leanok
Under the same hypotheses, fix a point of the original dynamic maximum. If
its whole maximum locus has empty trace in the endpoint over $U$, then every
point of the restricted actual distinct defect has full invariant strictly
smaller than the invariant at the fixed point.
\end{theorem}
\begin{proof}
Transport the local defect membership and its full value to the original
endpoint. The maximum property gives a weak inequality. If equality held,
the transported point would itself bound every full value on the original
defect and would therefore lie in the original maximum locus. It is also in
the endpoint over $U$, contradicting the empty-trace hypothesis. The weak
inequality is consequently strict.
\end{proof}
